\section{The main theorem}\label{subsec:main-theorem}\label{sec:theorem}

The theorem ties the continuum correlations to the original Wilson measures.
The lattice family is chosen once; the later choice of a volume subsequence determines one continuum hierarchy with all four conclusions below.

\begin{theorem}\label{thm:main}
For every gauge group $(m,G,\g)$ with $\g$ simple (Definition~\ref{def:gauge}), there exists a Wilson family $\W$ (Definition~\ref{def:family}) satisfying the regularity conditions~\ref{reg:1}--\ref{reg:4} of Definition~\ref{def:regular} such that, for every strictly increasing $p:\N\to\N$, there exist a strictly increasing $q:\N\to\N$ and a hierarchy over $\W$ (Definition~\ref{def:hierarchy}) with the following properties.
\begin{enumerate}[label=\textup{P\arabic*.},ref=P\arabic*]
\item\label{main:1} For every $n\in\N$, $v\in\mathcal T^n$ and separated test $f$ (Definition~\ref{def:tests}), the test is supported inside volume $p(q(k))$ eventually in $k$ (Definition~\ref{def:admissible}), and there exist $R_j\in\C$ such that
\begin{equation}\label{eq:main-limit}
\begin{aligned}
\lim_{N\to\infty}C_{N,j}(v;f)&=R_j
&&\text{($f$ supported in volume $j$)},\\
\lim_{k\to\infty}R_{p(q(k))}&=\int f(x)S^\mix(\pi(v_1),\ldots,\pi(v_n);x)\,dx\,.
\end{aligned}
\end{equation}
Here $C_{N,j}$ is the Wilson moment of Definition~\ref{def:moments}; these are the ordered limits of Definition~\ref{def:ordered}.
\item\label{main:2} The datum reconstructed from the selected point~$\xi_*$ satisfies the Wightman conditions~\ref{wightman:0}--\ref{wightman:3} of Definition~\ref{def:wightman}.
The hierarchy has all reconstruction properties~\ref{recon:1}--\ref{recon:15} of Definition~\ref{def:reconstructed}; in particular, its continuation has the selected phase's Euclidean correlations as values at ordered imaginary times.
The mixed correlations are the phase average
\begin{equation*}
S^\mix(a;x)=\int_X S^\xi(a;x)\,\nu(d\xi)\qquad(x\in\Conf n)\,.
\end{equation*}
\item\label{main:3} The single coefficient system in~\ref{hierarchy:7} satisfies both OPE estimates of Definition~\ref{def:ope}.
\item\label{main:4} The field $\Theta\in F_4$ has zero vacuum expectation as in~\ref{trace:1}, and its selected-phase two-point function (Definition~\ref{def:trace-pair}) satisfies
\begin{equation*}
\lim_{r\downarrow0}r^8\left(\log\frac{1}{r}\right)^2G_\Theta(r)
=\frac{3\dim_\R\g}{\pi^4}>0\,.
\end{equation*}
For every finite Wick profile $W$ (Definition~\ref{def:wick}) and every $\delta>0$, there exists $r\in(0,\delta)$ with $G_\Theta(r)\ne W(r)$.
These are the trace conditions~\ref{trace:1}--\ref{trace:4} of Definition~\ref{def:trace}.
\end{enumerate}
\end{theorem}

The lattice data are chosen before $p$.
The refinement, field spaces, phases, selected relativistic theory, product coefficients and trace field may depend on $p$.
The same refinement works for all orders of correlation, observable indices and separated tests.
The lattice limit in~\ref{main:1} is the mixture; the relativistic theory in~\ref{main:2} is attached to $\xi_*$.

The theorem does not choose a unique infinite-volume limit, show that different lattice choices give the same limit, or take the spacing and volume limits simultaneously.
It does not prove a mass gap, meaning a positive lower bound on nonzero energies above the vacuum.
It does not show that widely separated groups of observations become independent in the mixed state, that different phases agree, or that any one phase is itself the limit of lattice states.
The correlations are defined only when insertion points are distinct; the theorem does not define their smeared extension at coincident points.
The comparison excludes only the finite profiles in Definition~\ref{def:wick}; it does not exclude every polynomial combination of free fields or theories built from a continuum of free-field masses.
No Euclidean covariance or reflection-positivity assertion is included among~\ref{main:1}--\ref{main:4}; those properties are established internally in the construction.
